% ======================================================================
% PACKAGES
% ======================================================================
\usepackage[utf8]{inputenc}
\usepackage[T1]{fontenc}
\usepackage{lmodern}
\usepackage[english]{babel}
\usepackage{graphicx}
\usepackage{caption}
\usepackage{subfigure}
\usepackage[export]{adjustbox}
\usepackage{multirow}
\usepackage{booktabs}
\usepackage{colortbl}
\usepackage{xcolor}
\usepackage{hyperref}

\usepackage[backend=bibtex]{biblatex}
\addbibresource{references.bib}

% ======================================================================
% DESIGN -- Boadilla with a steel-blue primary color
% ======================================================================
\usetheme{Boadilla}
\captionsetup{font=scriptsize,labelfont=scriptsize}

% Primary color: standard SteelBlue (#4682B4).
% Structure controls headings, bullets, outline markers, and numbering.
\definecolor{steelblue}{HTML}{4682B4}
\setbeamercolor{structure}{fg=steelblue}
\setbeamercolor{titlelike}{fg=steelblue}
\setbeamercolor{block title}{fg=white,bg=steelblue}
\setbeamercolor{block body}{fg=black,bg=steelblue!8!white}
\hypersetup{colorlinks=true,linkcolor=steelblue,citecolor=steelblue,urlcolor=steelblue}

% Frame headings: section label above the bold slide title.
% The section label is automatic. An explicit \framesubtitle overrides it.
\setbeamerfont{frametitle}{series=\bfseries}
\setbeamerfont{framesubtitle}{size=\small,series=\mdseries}
\setbeamercolor{framesubtitle}{fg=steelblue}
\makeatletter
\defbeamertemplate{frametitle}{section above title}[1][left]{%
    \ifbeamercolorempty[bg]{frametitle}{}{\nointerlineskip}%
    \@tempdima=\textwidth%
    \advance\@tempdima by\beamer@leftmargin%
    \advance\@tempdima by\beamer@rightmargin%
    \begin{beamercolorbox}[sep=0.3cm,#1,wd=\the\@tempdima]{frametitle}
        \vbox{}\vskip-1ex%
        \if@tempswa\else\csname beamer@fte#1\endcsname\fi%
        {\usebeamerfont{framesubtitle}\usebeamercolor[fg]{framesubtitle}%
            \strut
            \ifx\insertframesubtitle\@empty
                \ifx\insertsectionhead\@empty
                    Overview%
                \else
                    \insertsectionhead
                \fi
            \else
                \insertframesubtitle
            \fi
            \strut\par}%
        \vskip0.1em%
        {\usebeamerfont{frametitle}\strut\insertframetitle\strut\par}%
        \vskip-1ex%
        \if@tempswa\else\vskip-.3cm\fi%
    \end{beamercolorbox}%
}
\setbeamertemplate{frametitle}[section above title]
\makeatother

% Left-align every title-page field, retaining the existing fonts and spacing.
\setbeamertemplate{title page}{%
    \vbox{}\vfill
    \begingroup
        \raggedright
        \begin{beamercolorbox}[sep=8pt,left]{title}
            \usebeamerfont{title}\inserttitle\par
            \ifx\insertsubtitle\empty\else
                \vskip0.25em
                {\usebeamerfont{subtitle}\usebeamercolor[fg]{subtitle}%
                    \insertsubtitle\par}%
            \fi
        \end{beamercolorbox}
        \vskip1em\par
        \begin{beamercolorbox}[sep=8pt,left]{author}
            \usebeamerfont{author}\insertauthor
        \end{beamercolorbox}
        \begin{beamercolorbox}[sep=8pt,left]{institute}
            \usebeamerfont{institute}\insertinstitute
        \end{beamercolorbox}
        \begin{beamercolorbox}[sep=8pt,left]{date}
            \usebeamerfont{date}\insertdate
        \end{beamercolorbox}
        {\usebeamercolor[fg]{titlegraphic}\inserttitlegraphic\par}
    \endgroup
    \vfill
}

% Outline slides use the default title layout, without a section label.
% The local group restores the section-above-title layout afterward.
\newcommand{\outlineframe}[1][]{%
    \begingroup
        \setbeamertemplate{frametitle}[default]
        \begin{frame}{Outline}
            \tableofcontents[#1]
        \end{frame}
    \endgroup
}

\AtBeginSection[]{%
    \outlineframe[currentsection]
}

\setbeamertemplate{footline}[frame number]
\setbeamertemplate{navigation symbols}{}
\setbeamertemplate{items}[square]
\setbeamertemplate{section in toc}[square]
\setbeamertemplate{itemize/enumerate body begin}{\normalsize}
\setbeamertemplate{itemize/enumerate subbody begin}{\normalsize}
\definecolor{light-gray}{gray}{0.9}

% ======================================================================
% HELPERS -- layout conveniences, without theme or font changes
% ======================================================================

% Unnumbered footnote, used for the title-page contact address.
\newcommand\blfootnote[1]{%
    \begingroup
    \renewcommand\thefootnote{}\footnote{#1}%
    \addtocounter{footnote}{-1}%
    \endgroup
}

% Usage: \figureplaceholder[height]{label}
% Replace the entire command with \includegraphics[width=\linewidth]{...}.
% The frame fits the current column width, including border and padding.
\newcommand{\figureplaceholder}[2][3cm]{%
    \fbox{%
        \parbox[c][#1][c]{\dimexpr\linewidth-2\fboxsep-2\fboxrule\relax}{%
            \centering #2\par
        }%
    }%
}
